The generation pipeline exhibits several failure modes—hallucinated or unsupported claims, missed or coarse citations, cross-section inconsistencies, JSON/schema and parsing errors, citation-to-BibTeX conflicts, provenance gaps, and failure cascades from poor retrieval—typically caused by incomplete retrieval, permissive prompts, decoding constraints, or heterogeneous metadata. These errors often interact and cascade, producing outputs that are syntactically valid but substantively incorrect. Conservative mitigations include requiring span-level citations and explicit source pointers, grounding checks with re-retrieval, stricter schema enforcement or post-hoc repairs, provenance augmentation, and human-in-the-loop review; related automated correction and detection work supports these directions \cite{web_thorne_2021_evidence_based_factual_error_correction,web_bayat_2023_fleek_factual_error_detection_and_correction_wit,web_macavaney_2018_overcoming_low_utility_facets_for_complex_answer}.

These strategies are intentionally conservative and their empirical effectiveness depends on retrieval quality and model capabilities; numeric error rates are reported in Results. Future priorities include controlled evaluations of mitigations, tighter integration of span-level factuality and claim-checking into the pipeline, and development of richer provenance formats to improve auditability and reproducibility.